% ============================================================
% Project : THE PHILOSOPHY OF VIABILITY
% Book    : Book I — Ontology of Viability
% Book ID : VS.B01
% File    : section_10_consistency_independence_and_conflict_analysis.tex
% Section : Chapter 07, Section 10
% Purpose : To establish criteria for examining consistency,
%           independence, incompatibility, and conflicts among
%           ontological commitments, definitions, axioms,
%           propositions, and formal interpretations.
% ============================================================

\section{Consistency, Independence, and Conflict Analysis}
\label{sec:consistency-independence-and-conflict-analysis}

\subsection{The Problem of Coherence in Formal Reconstruction}

A formal reconstruction must establish more than the intelligibility
of its individual propositions. It must also examine whether those
propositions can be maintained together, whether some depend on
assumptions that undermine others, and whether apparently distinct
claims introduce incompatible commitments.

The accumulation of formally expressed propositions does not
guarantee the coherence of the system to which they belong.
Definitions may introduce incompatible conditions. A proposition
may contradict an adopted commitment. Two independently plausible
claims may become incompatible when combined. A formal
representation may also create a conflict that does not belong
to the underlying ontological commitments but arises from an
inadequate choice of notation or interpretation.

The central task of this section is to establish a method for
distinguishing these possibilities.

Three questions must remain separate:

\begin{enumerate}
    \item \textbf{Consistency:} Can the relevant commitments and
    propositions be maintained together without contradiction
    under the specified formal and interpretive conditions?

    \item \textbf{Independence:} Can one claim be established or
    rejected without thereby establishing or rejecting another
    claim within the specified system?

    \item \textbf{Conflict:} Do two or more claims, definitions,
    interpretations, or commitments impose conditions that
    cannot jointly be satisfied?
\end{enumerate}

These questions are related, but they are not interchangeable.
A set of propositions may be consistent while containing
redundant claims. Two principles may be conceptually distinct
without being formally independent. A conflict between two
formal expressions may disappear when their different scopes
or meanings are clarified.

The governing thesis is that consistency, independence, and
conflict must be examined relative to explicitly identified
claims, assumptions, inference rules, and interpretations.
No conclusion about the coherence of the ontology should be
drawn from symbolic appearance alone.

\subsection{Consistency as a Relation within a Specified System}

Consistency is not adequately assessed by considering a
statement in isolation. It concerns the possibility of
maintaining a collection of statements together under specified
conditions.

Let \(\Gamma\) represent a set of formal premises and let
\(\mathcal{L}\) denote the relevant formal system. The set
\(\Gamma\) is syntactically consistent relative to
\(\mathcal{L}\) if it does not permit the derivation of a
contradiction.

A schematic formulation is:

\[
    \operatorname{Cons}_{\mathcal{L}}(\Gamma)
    \quad\text{iff}\quad
    \Gamma \nvdash_{\mathcal{L}} \bot,
\]

where \(\bot\) denotes contradiction.

This definition is relative to the adopted calculus. A different
formal system, with different inference rules or assumptions,
may produce a different result.

Semantic consistency is a related but distinct notion. A set
of premises is satisfiable when there is an admissible
interpretation or model in which all its members hold.

\[
    \operatorname{Sat}_{\mathcal{M}}(\Gamma)
    \quad\text{iff}\quad
    \exists \mathfrak{M}
    \bigl(\mathfrak{M}\models\Gamma\bigr).
\]

Here, \(\mathfrak{M}\) represents an interpretation satisfying
the formal requirements of the relevant language and
semantics.

The relationship between syntactic consistency and semantic
satisfiability depends on the properties of the adopted
formal system. Under an appropriate sound and complete
classical first-order framework, these notions are connected
through established metatheoretical results. Their equivalence
must not be presumed independently of the formal system.

This distinction is particularly important for an ontological
reconstruction because a formally consistent system may still
be ontologically inadequate. Its expressions may admit a
model while failing to represent the intended subject matter.

Formal consistency is therefore an important condition of
systematic adequacy, but it is not a complete criterion of
ontological truth.

\subsection{Different Levels of Consistency}

Consistency must be examined at more than one level because
different kinds of claims can conflict in different ways.

\paragraph{Formal consistency.}

Formal consistency concerns the compatibility of expressions
within a specified calculus. It asks whether the adopted
axioms, premises, and inference rules permit a contradiction
to be derived.

\paragraph{Semantic consistency.}

Semantic consistency concerns whether the expressions can
be jointly satisfied under an admissible interpretation.
It requires attention to the domain, predicates, relations,
and interpretive conditions of the formal language.

\paragraph{Conceptual consistency.}

Conceptual consistency concerns whether the meanings assigned
to the principal terms can be maintained together without
conceptual contradiction or equivocation.

A formal language may be syntactically well-formed while
its central terms are used inconsistently across different
arguments. Such a problem cannot be resolved merely by
checking derivations.

\paragraph{Ontological consistency.}

Ontological consistency concerns whether the commitments
and claims can be jointly maintained as claims about being
and identity.

A formal representation may be internally consistent while
its interpretation introduces ontological commitments that
conflict with the adopted foundations.

\paragraph{Methodological consistency.}

Methodological consistency concerns whether the procedures
used to formulate, compare, and justify claims are compatible
with the principles governing the inquiry.

For example, a methodological procedure that silently treats
recognition as a condition of existence would conflict with
the adopted ontological independence of existence from
cognition, even if the procedure were formally coherent
within a restricted model.

These levels must be examined separately before their
results are combined into an overall assessment.

\subsection{Consistency and the Three Foundational Principles}

The three foundational principles concerning being and
identity provide a central test case for consistency analysis.

Let:

\[
    AD = \text{Analytical Distinction},
\]

\[
    OI = \text{Ontological Inseparability},
\]

\[
    AP = \text{Absence of Ontological Priority}.
\]

Their joint adoption is represented as:

\[
    F_{BI} := AD \land OI \land AP.
\]

The conjunction expresses their simultaneous adoption.
It does not establish their formal consistency automatically.
Their compatibility must be examined through the meanings
assigned to each principle and through the adequacy of their
formal representations.

\paragraph{Analytical Distinction and Ontological Inseparability.}

At first sight, these principles may appear to conflict.
If being and identity are analytically distinct, how can
they also be ontologically inseparable?

The apparent conflict arises only if analytical distinction
is interpreted as ontological separation.

Analytical distinction concerns the difference between
the questions or roles through which being and identity
are examined. Ontological inseparability concerns their
relation within the ontology.

The two principles therefore address different aspects
of the inquiry.

Their joint adoption requires the reconstruction to preserve
both claims:

\begin{itemize}
    \item being and identity must remain analytically
    distinguishable;
    \item their analytical distinction must not be converted
    into an assertion of independent ontological existence.
\end{itemize}

No contradiction follows merely from maintaining these
two commitments together. However, their compatibility
depends on a sufficiently precise account of what analytical
distinction and ontological inseparability mean.

\paragraph{Ontological Inseparability and Absence of Ontological
Priority.}

These principles are also compatible in their intended
formulations.

Ontological inseparability concerns the absence of ontological
separability between being and identity. Absence of ontological
priority denies that either is ontologically prior to the
other in the specified sense.

The two principles are related but not equivalent. Inseparability
does not automatically establish the absence of every form
of priority, and the denial of priority does not independently
establish inseparability.

Their joint adoption is therefore a substantive combination
of distinct commitments rather than a repetition of one
principle in two formulations.

\paragraph{Analytical Distinction and Absence of Ontological
Priority.}

Analytical distinction does not establish that one analytical
question is ontologically prior to the other. The distinction
between the questions is compatible with the denial of
ontological priority.

Likewise, the denial of ontological priority does not erase
the analytical distinction.

The principles can therefore be maintained together, provided
that analytical ordering is not confused with ontological
priority.

\paragraph{Joint assessment.}

The three principles are conceptually compatible under their
intended meanings. This assessment is not yet a formal proof
of consistency for every possible formalization of them.

A formal consistency result would require a sufficiently
specified calculus, an explicit interpretation of the
relevant symbols, and an examination of the resulting
system.

The present reconstruction therefore adopts their compatibility
as a coherent foundational orientation while retaining
the formal demonstration of consistency as a technical
task requiring further specification.

\subsection{Apparent Contradiction and Genuine Contradiction}

A contradiction must be distinguished from a difference
of emphasis, level, scope, or terminology.

Two claims are genuinely contradictory only when their
content and conditions make it impossible to maintain
both under the same relevant interpretation.

For example, in classical logic, the conjunction

\[
    \varphi \land \neg\varphi
\]

constitutes a contradiction.

However, two apparently opposed statements may not have
this structure.

Consider the following pair:

\begin{quote}
Being and identity are analytically distinct.

Being and identity are ontologically inseparable.
\end{quote}

These statements are not contradictory when their predicates
refer to different aspects of the relation. The first concerns
analytical differentiation; the second concerns ontological
inseparability.

A contradiction would arise only if the first statement
were reformulated to assert ontological separation in the
same respect in which the second asserts inseparability.

The evaluation of an apparent conflict must therefore
examine:

\begin{enumerate}
    \item the meanings of the terms involved;
    \item the level of analysis at which each claim operates;
    \item the domain and scope of each claim;
    \item the respect in which the claims are asserted;
    \item the temporal or conditional qualifications, if any;
    \item the interpretation under which their joint
    satisfaction is being evaluated.
\end{enumerate}

Without this examination, the reconstruction may mistake
a legitimate analytical distinction for an inconsistency,
or conceal a genuine contradiction by using different
terminology for incompatible commitments.

\subsection{Local and Global Consistency}

A set of claims may be consistent when considered separately
but inconsistent when combined with other claims.

This distinction requires the examination of local and
global consistency.

\paragraph{Local consistency.}

Local consistency concerns a restricted group of claims,
such as the propositions within one section or the formal
representation of one foundational principle.

A section may be locally coherent while depending on an
assumption that conflicts with a commitment established
elsewhere.

\paragraph{Global consistency.}

Global consistency concerns the relevant collection of
claims across the complete formal reconstruction.

It requires examining the interaction of foundations,
definitions, propositions, interpretations, and additional
assumptions throughout the system.

Local consistency does not guarantee global consistency.
A proposition may be valid within its immediate argument
while becoming incompatible with a claim introduced in
another section.

The reconstruction must therefore combine local examination
with periodic global audits.

This is especially important when the same concept appears
in several chapters. A definition of identity introduced
in one context may acquire additional conditions in another.
Those conditions must be examined to determine whether they
refine, extend, or contradict the earlier formulation.

\subsection{Independence: Conceptual and Formal Dimensions}

Independence concerns whether one claim can be maintained
without the other, but the term has several meanings that
must not be conflated.

\paragraph{Conceptual independence.}

Two claims are conceptually independent when they perform
distinct conceptual functions and cannot be reduced to
one another merely by clarifying their meanings.

Conceptual independence does not imply that the claims
have no relation. Two principles may be mutually relevant
while retaining distinct conceptual content.

\paragraph{Logical independence.}

Two expressions are logically independent relative to
a specified formal system when neither is derivable from
the other under the adopted premises and rules, subject
to the relevant definition of independence.

\paragraph{Axiomatic independence.}

An axiom \(\alpha\) is independent of a set of other axioms
\(\Delta\) when it cannot be derived from \(\Delta\) within
the specified formal system.

A standard model-theoretic method for demonstrating such
independence is to construct an admissible interpretation
in which the axioms in \(\Delta\) hold while \(\alpha\)
does not.

\paragraph{Justificatory independence.}

Two claims are justificatorily independent when the grounds
for accepting one do not rely on the acceptance of the
other, even if the claims are conceptually related.

These dimensions must be distinguished. Conceptual difference
does not prove formal independence. The absence of an
identified derivation does not establish independence.
Justificatory independence requires examination of the
actual grounds supporting each claim.

\subsection{Independence of the Three Foundational Principles}

The three principles concerning being and identity must
retain their distinct roles.

Analytical Distinction establishes a difference between
the analytical questions or roles associated with being
and identity.

Ontological Inseparability establishes that their relation
cannot be adequately represented as ontological separation.

Absence of Ontological Priority denies the specified
priority relation in either direction.

Their conceptual distinction is necessary to preserve
the architecture of the ontology. If one principle were
simply replaced by another, the reconstruction would lose
a substantive distinction.

However, conceptual distinction does not by itself prove
that the formal representations of the three principles
are independent within every possible calculus.

A formal independence result would require:

\begin{enumerate}
    \item an explicitly specified formal language;
    \item a defined set of axioms and inference rules;
    \item precise formal representations of the principles;
    \item an appropriate semantic framework;
    \item models or other demonstrations establishing
    the required independence relations.
\end{enumerate}

Until these conditions are satisfied, the reconstruction
should describe the principles as conceptually and
justificatorily distinct foundational commitments, while
leaving their full formal independence as an open technical
question.

\subsection{The Limits of Formal Equivalence}

Two claims may be equivalent within a particular formal
system without being conceptually identical in every
relevant sense.

If:

\[
    \Gamma \vdash_{\mathcal{L}} (\varphi \leftrightarrow \psi),
\]

then \(\varphi\) and \(\psi\) are interderivable relative
to \(\Gamma\) and \(\mathcal{L}\), assuming the adopted
interpretation of the biconditional.

This does not automatically establish that the concepts
represented by \(\varphi\) and \(\psi\) are identical in
their full ontological significance.

Formal equivalence may depend on restricted definitions,
a particular domain, or assumptions that do not exhaust
the meanings of the relevant concepts.

This distinction is especially important for the relation
between being and identity.

The co-applicability schema:

\[
    \forall x \in D_E
    \bigl(E(x) \rightarrow (B(x) \land I(x))\bigr)
\]

does not establish:

\[
    \forall x \in D_E
    \bigl(B(x) \leftrightarrow I(x)\bigr).
\]

The first expression concerns the joint applicability
of two predicates under specified conditions. The second
expresses a biconditional relation between their conditions
of application.

Even if a particular formalization were to establish a
biconditional under restricted assumptions, further
argument would still be required before treating it as
an adequate representation of ontological identity.

The reconstruction must therefore distinguish formal
equivalence, conceptual equivalence, and ontological
identity.

\subsection{Types of Conflict}

Conflicts within the reconstruction may arise in different
ways. Identifying their type is necessary before deciding
whether a claim must be rejected, revised, restricted,
or reinterpreted.

\paragraph{Logical conflict.}

A logical conflict occurs when a set of propositions
entails a contradiction under the adopted formal system.

For example:

\[
    \Gamma \vdash_{\mathcal{L}} \varphi
\]

and

\[
    \Gamma \vdash_{\mathcal{L}} \neg\varphi.
\]

The resulting conflict concerns the inferential consequences
of the premises and rules.

\paragraph{Semantic conflict.}

A semantic conflict occurs when the meanings assigned
to expressions cannot be jointly satisfied under the
specified interpretation.

Such a conflict may arise even when the expressions
are individually well-formed.

\paragraph{Conceptual conflict.}

A conceptual conflict occurs when a term is assigned
incompatible meanings or conditions of application
within the same theoretical context.

It may result from equivocation, an unstable definition,
or the unnoticed transfer of a concept between different
levels of analysis.

\paragraph{Ontological conflict.}

An ontological conflict occurs when two adopted claims
impose incompatible conditions concerning being, identity,
or their relations.

Such a conflict cannot be resolved merely by changing
notation. It requires reconsideration of the substantive
ontological commitments.

\paragraph{Methodological conflict.}

A methodological conflict occurs when the procedures
used to justify claims contradict the adopted standards
of inquiry.

For example, a procedure may treat a conclusion as
established while relying on an assumption that the
methodology classifies as provisional.

\paragraph{Scope conflict.}

A scope conflict occurs when claims are applied to
incompatible domains or levels without an explicit
justification for their extension.

A statement that is valid under a restricted interpretation
may become incompatible with another statement when
incorrectly generalized.

\paragraph{Terminological conflict.}

A terminological conflict occurs when different expressions
are used for the same intended concept without clarification,
or when the same expression is used for different concepts
without acknowledging the change.

This kind of conflict can create the appearance of
theoretical disagreement or conceal an actual disagreement.

\subsection{A Method for Diagnosing Apparent Conflicts}

An apparent conflict should not immediately lead to the
rejection of one of the claims. The relation between the
claims must first be reconstructed.

The following procedure is required:

\begin{enumerate}
    \item \textbf{Identify the claims.}
    State the two or more claims whose compatibility is
    in question.

    \item \textbf{Clarify their meanings.}
    Determine whether the principal terms are used in
    the same sense.

    \item \textbf{Specify their scopes.}
    Identify the domains, conditions, and respects in
    which each claim is asserted.

    \item \textbf{Identify their levels.}
    Determine whether the claims concern ontology,
    conceptual articulation, formal representation,
    methodology, or interpretation.

    \item \textbf{Reconstruct their commitments.}
    Make explicit the premises and assumptions required
    by each claim.

    \item \textbf{Test joint satisfiability.}
    Examine whether both claims can hold under a common
    admissible interpretation.

    \item \textbf{Locate the source of conflict.}
    Determine whether the problem is logical, semantic,
    conceptual, ontological, methodological, or related
    to scope.

    \item \textbf{Determine the appropriate response.}
    Decide whether the claims require clarification,
    qualification, reinterpretation, restriction,
    reformulation, or rejection.
\end{enumerate}

This procedure prevents the premature elimination of
claims that may be compatible once their meanings and
conditions are specified.

It also prevents the opposite error: treating a genuine
contradiction as a harmless difference of terminology.

\subsection{Contradiction, Tension, and Theoretical Incompleteness}

Not every unresolved relation between claims constitutes
a contradiction.

A theoretical tension exists when two claims are difficult
to reconcile or appear to pull the theory in different
directions, but their incompatibility has not been
established.

For example, the commitment to ontological inseparability
and the requirement of analytical distinction create
a demanding explanatory task. The ontology must explain
how two aspects can be distinguished without being
ontologically separated.

This is a theoretical tension in the sense that it
requires a sufficiently precise account. It is not,
without further argument, a contradiction.

Theoretical incompleteness occurs when the available
definitions, arguments, or formal representations do
not yet determine whether two claims are compatible.

Incompleteness should not be confused with inconsistency.
A theory may be incomplete because it lacks the resources
to settle a question, while remaining consistent within
the claims it has actually adopted.

The reconstruction must therefore preserve at least
three distinct diagnoses:

\begin{itemize}
    \item \textbf{Contradiction:} the claims cannot jointly
    hold under the specified conditions.

    \item \textbf{Theoretical tension:} the claims require
    further explanation or reconciliation, but incompatibility
    has not been established.

    \item \textbf{Incompleteness:} the available framework
    does not yet determine the relation between the claims.
\end{itemize}

These distinctions are essential for avoiding both
premature rejection and unjustified reconciliation.

\subsection{The Effect of Inconsistency on a Formal System}

In classical logic, an inconsistent set of premises may
permit the derivation of arbitrary conclusions through
the principle commonly known as explosion.

In schematic form:

\[
    \varphi,\neg\varphi \vdash_{\mathcal{L}} \psi.
\]

If the adopted calculus permits this inference, an
uncontrolled contradiction can undermine the usefulness
of the system as a means of distinguishing justified
consequences from unrelated claims.

This problem makes the identification and treatment
of contradictions especially important.

However, the existence of an inconsistency in one part
of a theoretical reconstruction does not necessarily
mean that every claim in the broader philosophical
project must be abandoned. The scope of the inconsistency,
the rules governing inference, and the dependencies
of the affected claims must be examined.

Some non-classical logics are designed to manage
inconsistent information without permitting arbitrary
conclusions. Their availability does not mean that the
present ontology has adopted a paraconsistent logic.
Such a decision would require a separate justification
and an explicit account of the consequences for the
formal reconstruction.

The present section therefore retains the requirement
that contradictions be identified and resolved wherever
possible, without prematurely selecting a non-classical
logic as a substitute for conceptual clarification.

\subsection{Consistency Under Alternative Formalizations}

A foundational commitment may admit more than one formal
representation.

Different formalizations may preserve the same intended
ontological content while differing in their expressive
resources, inferential rules, or semantic assumptions.

Consequently, the consistency of one formalization does
not automatically establish the consistency of every
alternative representation.

Likewise, an inconsistency in one formalization does
not automatically demonstrate that the underlying
ontological commitment is itself contradictory.

The evaluation must distinguish:

\begin{enumerate}
    \item the consistency of the foundational commitment
    as conceptually articulated;
    \item the consistency of a particular formal representation;
    \item the adequacy of the interpretation assigned to
    that representation;
    \item the compatibility of the representation with
    other adopted commitments.
\end{enumerate}

If two formalizations produce different results, the
difference must be examined to determine whether it
arises from their logical systems, their definitions,
their interpretations, or their actual ontological
commitments.

The formal system must not be treated as an unquestionable
measure of the ontology. Its adequacy remains subject
to the philosophical distinctions and constraints
that motivated its construction.

\subsection{Consistency and Independence Testing Through Models}

Models can provide a method for examining consistency
and independence, provided that the formal language
and its semantics are sufficiently specified.

A model \(\mathfrak{M}\) satisfies a set of expressions
\(\Gamma\) when every member of \(\Gamma\) is true
under the interpretation defined by \(\mathfrak{M}\).

If such a model exists, the set is satisfiable under
the relevant semantics:

\[
    \exists \mathfrak{M}
    \bigl(\mathfrak{M}\models\Gamma\bigr).
\]

To examine the independence of an expression \(\alpha\)
from a set \(\Delta\), one possible method is to construct
a model in which:

\[
    \mathfrak{M}\models\Delta
\]

while:

\[
    \mathfrak{M}\not\models\alpha.
\]

Under suitable conditions, such a model demonstrates
that \(\alpha\) is not a semantic consequence of
\(\Delta\).

A corresponding test may be conducted for the negation
of \(\alpha\), depending on the precise independence
claim being investigated.

These procedures require caution. A model may demonstrate
a formal property of a representation without demonstrating
that the model corresponds to an actual ontological
possibility.

The construction of a model therefore establishes
a result relative to the specified semantics. Its
ontological interpretation requires additional argument.

For the present reconstruction, model-based testing
should be used to examine formal consistency and
independence only after the relevant language, domains,
predicates, relations, and axioms have been adequately
specified.

\subsection{A Consistency and Conflict Register}

The reconstruction requires a systematic record of
the results of consistency and conflict analysis.

For every substantial claim or group of claims, the
following information should be recorded:

\begin{enumerate}
    \item \textbf{Claim identification:} the exact proposition,
    commitment, definition, or formal expression under
    examination.

    \item \textbf{Functional form:} its classification as
    a foundational commitment, definition, axiom,
    proposition, thesis, hypothesis, argument, or
    open problem.

    \item \textbf{Epistemic status:} the grounds and
    degree of acceptance assigned to it.

    \item \textbf{Scope:} the domain, conditions, and
    respect in which it applies.

    \item \textbf{Dependencies:} the premises, definitions,
    interpretations, and rules required for its support.

    \item \textbf{Consistency result:} whether it is
    consistent with the specified set of other claims.

    \item \textbf{Independence result:} whether any
    relevant conceptual, logical, axiomatic, or
    justificatory independence has been established.

    \item \textbf{Conflict classification:} whether an
    apparent conflict is logical, semantic, conceptual,
    ontological, methodological, terminological, or
    scope-related.

    \item \textbf{Resolution:} whether the issue has
    been clarified, qualified, reformulated, restricted,
    rejected, or retained as unresolved.

    \item \textbf{Revision consequences:} the other
    claims that must be reconsidered if the examined
    claim is modified.
\end{enumerate}

The register should distinguish demonstrated results
from provisional assessments. A claim should not be
marked formally consistent merely because no
contradiction has yet been identified.

Likewise, a claim should not be marked independent
merely because its derivation from another claim has
not yet been found.

The register is an instrument of intellectual control,
not a substitute for the arguments through which
consistency and independence are established.

\subsection{Open Problems}

Several issues remain unresolved at the present stage
of the formal reconstruction.

First, the complete formal language and calculus have
not yet been fixed. Consequently, a definitive global
consistency proof for the entire formal system cannot
yet be supplied.

Second, the formal independence of the three foundational
principles remains to be investigated within a sufficiently
specified system. Their conceptual and justificatory
distinction is established as a requirement of the
ontology, but their technical independence requires
further work.

Third, the complete formal representation of ontological
inseparability remains open. The co-applicability schema
does not exhaust the intended ontological content, and
a stronger representation must avoid collapsing
analytical distinction into predicate equivalence.

Fourth, the conditions under which different formal
representations preserve the same ontological commitments
require further examination.

Fifth, a systematic global conflict audit must be
performed across the complete reconstruction. Local
consistency within individual sections cannot substitute
for this broader examination.

Sixth, the relation between formal consistency and
ontological adequacy requires continued clarification.
A formally satisfiable system may still fail to represent
the intended ontological structure.

These questions remain open without implying that the
foundational architecture is already inconsistent.
They identify the additional demonstrations required
before the formal reconstruction can be regarded as
technically complete.

\subsection{Concluding Principle}

Consistency, independence, and conflict analysis are
necessary because a coherent ontology must preserve
the compatibility of its commitments while allowing
their conceptual differences to remain explicit.

Consistency requires that claims can be maintained
together under specified conditions. Independence
requires a demonstration appropriate to the dimension
being examined. Conflict analysis requires distinguishing
genuine incompatibility from differences of scope,
meaning, level, or interpretation.

The governing principle may therefore be stated as follows:

\begin{quote}
No claim should be accepted, rejected, or identified
as independent solely on the basis of its formal
appearance or conceptual description. Its compatibility
with other claims, its inferential relations, and
the conditions of its interpretation must be examined
within an explicitly specified framework.
\end{quote}

This principle provides a basis for testing the formal
reconstruction without confusing conceptual distinction
with contradiction, formal consistency with ontological
truth, or the absence of an established derivation
with demonstrated independence.

The next section examines the role of models,
countermodels, and alternative interpretations in
testing the consequences and limitations of the
formal system.

% ============================================================
% END OF SECTION 10 OF CHAPTER 07
% ============================================================